% Statement and cited-input packet. Target proof is separately embargoed.
% Frozen II source SHA256: 74497bfe327fae25c9cdcb17334b3bf42a9793a9df7c0a4404d97568f2f760f0
% Target owner: cdp:st:D.2.8; tag: 01GF

\begin{proposition}[Proved coherent and integral comparisons for \ref{cdp:st:D.2.7}]\hypertarget{tag-01GF}{}\label{cdp:st:D.2.8}



For the exact threefold and markings of \ref{cdp:st:D.2.7},
\[
 R^0f_*\OO_Y=\OO_{\PP^1},\qquad
 R^1f_*\OO_Y=\OO_{\PP^1}^{\oplus2},\qquad
 R^2f_*\OO_Y=\OO_{\PP^1},
 \qquad h^q(Y,\OO_Y)=(1,2,1,0).
\tag{D.33}\hypertarget{tag-019K}{}\label{cdp:eq:D.33}
\]
The coherent restriction map to every full fibre is surjective.
The integral comparison is
\[
 H^0(\PP^1,R^1f_*\Z)
   =\Z(e_1^*+e_2^*-e_4^*)
   \ \subsetneq\ H^1(A;\Z)^{\pi_1}=\Z^4.
\tag{D.34}\hypertarget{tag-019L}{}\label{cdp:eq:D.34}
\]
The map in \eqref{cdp:eq:D.34} is restriction to the smooth base, and the
display gives its exact image. The exceptional torsion locus
for restrictions of $\Pic(Y)=H^1(Y,\OO_Y)\simeq\C^2$ to
at least one fibre is exactly
\[
 \iota\bigl(H^1(A;\Q)\bigr)\subset H^1(A,\OO_A),
 \qquad
 \iota(a,b,c,d)=\tfrac12(a+\mathrm i b,c+\mathrm i d).
\tag{D.35}\hypertarget{tag-019M}{}\label{cdp:eq:D.35}
\]
It is countable and dense. Thus the all-fibre non-torsion locus
of \ref{cdp:st:D.1.11} need not be open, even under integral
$H^1(Y;\Z)=H^2(Y;\Z)=0$.
\end{proposition}

% Cited input cdp:st:D.1.11; statement or labelled formula only.
\begin{proposition}\hypertarget{tag-01FW}{}\label{cdp:st:D.1.11}


Let $f:Y\to\PP^1$ be a proper connected-fibre holomorphic map from
a connected compact complex manifold with $b_1(Y)=b_2(Y)=0$.
In every component of $\Pic(Y)$, the line
bundles which are non-torsion on every full fibre form a dense
$G_\delta$ set of full Lebesgue measure in its complex vector-space
coordinates. If evaluation of global sections of $R^1f_*\OO_Y$
has rank one, then $\Pic^0(Y)=\C$ and the exceptional set is
countable. For two-torus fibres this alternative holds whenever
the splitting has a negative summand.
In particular the general-line-bundle conclusion of CDP20
Proposition 5.6 holds with this precise quantifier, without the
nonnegative-summand assertion of Proposition 5.5.
\end{proposition}

% Cited input cdp:st:D.1.12; statement or labelled formula only.
\begin{proposition}\hypertarget{tag-01FX}{}\label{cdp:st:D.1.12}


For a proper connected-fibre holomorphic map $f:Y\to C$ from a
compact complex threefold with $b_2(Y)=0$ and general fibre a
two-torus or a primary Kodaira surface, the cup-product homomorphism
is an isomorphism
\[
\bigwedge^2R^1f_*\OO_Y\xrightarrow{\simeq}R^2f_*\OO_Y.
\]
Every full fibre has structure-sheaf cohomology dimensions
$(h^0,h^1,h^2)=(1,2,1)$, and these two direct images satisfy
base change at every fibre. This proves these two cases of CDP20
Proposition 5.1 by an index and cup-product argument.
\end{proposition}

% Cited input cdp:st:D.1.7; statement or labelled formula only.
\begin{proposition}\hypertarget{tag-01G5}{}\label{cdp:st:D.1.7}


For a connected compact complex space $S$, including a nonreduced
one, the coefficient map
\[
H^1(S;\mathbb R)\longrightarrow H^1(S,\OO_S)
\]
is injective as a map of real vector spaces.
\end{proposition}

% Cited input cdp:st:D.2.7; statement or labelled formula only.
\begin{proposition}[Proved existence]\hypertarget{tag-01GE}{}\label{cdp:st:D.2.7}



There exists a smooth compact connected complex threefold $Y$
with a holomorphic map $f:Y\to\PP^1$ with connected fibres such
that
\[
 H_1(Y;\Z)=H_2(Y;\Z)=0,\qquad
 a(Y)=1,\qquad c_3(Y)=0,\qquad b_3(Y)=2.
\tag{D.21}\hypertarget{tag-0198}{}\label{cdp:eq:D.21}
\]
Every smooth fibre is the fixed abelian surface
$A=E\times E$, where $E=\C/(\Z+\mathrm i\Z)$, and the marked
integral monodromy on $H_1(A;\Z)$ is the identity. There are
exactly three singular fibres. Each is reduced with simple
normal crossings, and consists of a primary Kodaira surface
and a primary Hopf surface meeting along one smooth elliptic
curve. Consequently the trivial-monodromy class in D.2 is
nonempty, and equality in the unrestricted sign theorem
\ref{cdp:st:D.2.5} is attained.
\end{proposition}

% Cited input cdp:st:D.2.5; statement or labelled formula only.
\begin{proposition}[Proved unrestricted sign under trivial monodromy]\hypertarget{tag-01GC}{}\label{cdp:st:D.2.5}



Let $f:Y\to\PP^1$ be a proper surjective holomorphic map with
connected fibres from a compact connected complex threefold with
$b_1(Y)=b_2(Y)=0$. Suppose its smooth general fibre is a complex
two-torus and its monodromy representation on $H_1(F,\Q)$ is
trivial. Then
\[
 c_3(Y)\le0,\qquad b_3(Y)\ge2.
 \tag{D.S.1}\hypertarget{tag-019O}{}\label{cdp:eq:D.S.1}
\]
The fibres may have arbitrary singularities and arbitrary
multiplicities. In particular this proves the sign assertion of
D.2 for a holomorphic algebraic reduction; no existence assertion
is required for the conclusion.
\end{proposition}

